%% ==========================================================================
%% SECTION 4: THE IMPOSSIBILITY
%% ==========================================================================
\section{The Impossibility}\label{sec:formal}

% The previous section established that internal epistemic state exists and is
% discarded by the text interface. The system is unobservable in precisely the
% dimension that matters for verification. We now formalize this observation as an
% impossibility result: under text-only observation, no bounded supervisor can
% distinguish honest systems from dishonest ones.

A common assumption in current practice is that epistemic honesty is a
capability problem: sufficiently capable models, trained with sufficient RLHF,
will eventually ``know what they know'' and report it accurately. Scaling solved
grammar; scaling solved coherence; scaling will solve fabrication. Lin et
al.\ tested this assumption directly and found an inverse scaling trend: larger
models were \emph{less} truthful on TruthfulQA~\cite{lin2022truthfulqa}. The
theorems that follow explain why this expectation is architecturally
incoherent. The constraint is not capability but interface or, in cost terms,
the wrong investment tier. Text-only observation lacks the
degrees of freedom to verify epistemic state, regardless of what the model
internally represents. No amount of capability improvement can escape a
verification impossibility because the model may know perfectly well that it is
fabricating, but the interface provides no channel through which this knowledge
can be verified.

Our argument proceeds in two stages. First, we show that any policy conditioning
only on the query cannot satisfy epistemic honesty across ambiguous world states
(\autoref{thm:representational}: Representational Impossibility). Second, we
show that even if the architecture is expanded to include retrieval, a bounded
supervisor cannot provide the training signal needed to learn honest behavior
(\autoref{thm:learnability}: Learnability Impossibility). Together, these
results establish that epistemic honesty is unachievable for text-only
observation models under realistic supervision constraints.

To be precise about the type of claim: this is a \emph{verification impossibility}.
We prove that a bounded supervisor observing only text cannot distinguish
epistemically honest systems from systems that fabricate plausibly. The
impossibility is architectural: it applies to any system operating under text-only
observation, regardless of model family, scale, or training procedure. These
results are prompt-agnostic: they hold for any policy that exports only text,
regardless of whether a system prompt instructs the model to be ``helpful,''
``honest,'' or otherwise.

\subsection{Definitions}

Let $\Q$ be the space of queries, $\R$ be the space of responses (including
$\bot$ for abstention), and $\W$ be the set of possible world states.

\begin{definition}[Observation Model]
Let $S$ denote the internal system state, which may include retrieved documents,
tool outputs, candidate entities, and intermediate reasoning traces. An
observation model is a function $O : S \rightarrow \Omega$ that determines what
information is exposed to the supervisor or user.

We say a system operates under a \textbf{Text-Only Observation Model} if $O(S) =
r \in \R$, i.e., the exported observation consists solely of the linearized
response string (optionally including token-level probabilities), and excludes
any structured epistemic provenance or causal trace.
\end{definition}

\begin{definition}[Predictor-Centric Policy]
For any query $q \in \Q$, $r \in \R$, and world state $w \in \W$, a
predictor-centric model is defined by a stochastic policy $\pi: \Q \to
\Delta(\R)$, where $\pi(r|q)$ denotes the probability of generating response $r$
given query $q$. Crucially, $\pi$ does not take $w$ as an input at inference
time.
\end{definition}

\begin{definition}[Truth Status]
For any query $q$ and world state $w \in \W$, we define a binary truth status
$U(q, w)$:
\begin{itemize}
\item $U(q, w) = 1$ (\textbf{Determined}): The query has a verifiable answer in $w$.
\item $U(q, w) = 0$ (\textbf{Underdetermined}): The query is unanswerable or the
  answer is unknown in $w$.
\end{itemize}
\end{definition}

\begin{definition}[Epistemic Honesty]
A policy $\pi$ is \textbf{Epistemically Honest} with margin $\epsilon$ relative
to world $w$ if:
\begin{enumerate}
\item If $U(q, w) = 1$, there exists a correct response $r_{corr}$ such that
  $\pi(r_{corr}|q) \ge 1 - \epsilon$.
\item If $U(q, w) = 0$, $\pi(\bot|q) \ge 1 - \epsilon$.
\end{enumerate}
\end{definition}

\subsection{Representational Impossibility}

\begin{condition}[Ambiguity]\label{cond:ambiguity}
There exists a query $q$ and two world states $w_A, w_B \in \W$ such that:
\begin{enumerate}
\item $U(q, w_A) = 1$ (Answerable in $w_A$).
\item $U(q, w_B) = 0$ (Unanswerable in $w_B$).
\end{enumerate}
\end{condition}

\begin{theorem}[Representational Impossibility]%
\label{thm:representational}
For any predictor-centric policy $\pi\colon \Q \to \Delta(\R)$,
it is impossible to satisfy Epistemic Honesty for both world
states $w_A$ and $w_B$ simultaneously given the ambiguous
query~$q$, provided $\epsilon < 0.5$.
\end{theorem}

\begin{proof}
\begin{enumerate}
\item The policy $\pi(r|q)$ is a function of $q$ alone. Thus, the distribution
  over responses is identical in both worlds: $\pi(\cdot|q, w_A) = \pi(\cdot|q,
  w_B) = \pi(\cdot|q)$.
\item To be honest in $w_A$, we require $\pi(r_{corr}|q) \ge 1 - \epsilon$.
\item To be honest in $w_B$, we require $\pi(\bot|q) \ge 1 - \epsilon$.
\item Since $r_{corr} \neq \bot$, these events are disjoint. The sum of
  probabilities would be $\ge 2(1-\epsilon)$.
\item If $\epsilon < 0.5$, the sum exceeds 1, which violates the axioms of
  probability.
\item Therefore, no such policy $\pi$ exists. \qedhere
\end{enumerate}
\end{proof}

\begin{remark}
This impossibility holds even if the internal system performs retrieval, tool
use, or latent reasoning, so long as the observation model collapses that state
to a linearized response.
\end{remark}

\fakepara{A concrete example.}
Consider the query: ``What did Dr.\ Yamamoto's 2021 study conclude
about mitochondrial decay?''

In world $w_A$, Dr.\ Yamamoto exists and published such a study. The honest
response is the study's conclusion. In world $w_B$, no such researcher or study
exists. The honest response is abstention.

The policy sees only the query $q$. It cannot distinguish $w_A$ from $w_B$. Yet
honesty requires answering in $w_A$ and abstaining in $w_B$. The same query, two
incompatible requirements, one policy. This is the impossibility in miniature:
the interface lacks the degrees of freedom to express what honesty requires.

\fakepara{Connection to identifiability.}
In statistical terms, \autoref{thm:representational} states that epistemic
honesty is \emph{non-identifiable} under text-only observation: the mapping from
world states to observable distributions is many-to-one. Worlds $w_A$ and $w_B$
induce identical observable distributions $\pi(\cdot|q)$ despite requiring
different honest behaviors. No amount of data from this observation channel can
distinguish them. This connects the impossibility to classical results on
observational equivalence in statistics and diagnosability in systems theory.

\subsection{The RAG Escape?}

\autoref{thm:representational} might seem to have an obvious escape: give
the model access to external information. Retrieval-augmented generation does
exactly this: the policy conditions on retrieved documents, not just the query.
Does this break the impossibility?

It does not. The problem shifts from representation to learning.

\begin{definition}[Grounded Policy]
A grounded policy is a stochastic function $\pi: \Q \times \D \to \Delta(\R)$,
where $\D$ represents retrieved documents or observations of $w$ available at
inference time. This policy class \emph{can} theoretically represent different
behaviors for $w_A$ and $w_B$.
\end{definition}

\begin{definition}[Bounded Supervisor]
The training process relies on a supervisor $S$ that assigns reward $R$ based on
an observation function $Obs(q, r, w)$. The supervisor has a verification budget
$B$.
\[
Obs(q, r, w) =
\begin{cases}
(q, r, \verify(r, w)) & \text{if } C_{verify}(q, r, w) \le B \\
(q, r, \emptyset) & \text{if } C_{verify}(q, r, w) > B
\end{cases}
\]
The supervisor's policy may be human, automated, or hybrid, but is constrained
to act solely on the observation $Obs(q,r,w)$ and is subject to a finite
verification budget $B$.
\end{definition}

\begin{condition}[Hallucination Regime]\label{cond:hallucination}
There exists a plausible fabrication $r_{fab}$ and contexts $(q, w_A), (q, w_B)$
from \autoref{cond:ambiguity} such that:
\begin{enumerate}
\item The cost to verify $r_{fab}$ in both contexts exceeds $B$ (e.g., obscure
  or non-existent citation).
\item Consequently, the supervisor observation is identical:
  \[ Obs(q, r_{fab}, w_A) = Obs(q, r_{fab}, w_B) = (q, r_{fab}, \emptyset) \]
\end{enumerate}
\end{condition}

\begin{theorem}[Learnability Impossibility]\label{thm:learnability}
No learning algorithm $\mathcal{A}$ optimizing reward $R$ from a bounded
supervisor $S$ can converge to a Grounded Policy that satisfies Epistemic
Honesty across both contexts $w_A$ and $w_B$ for the response $r_{fab}$.
\end{theorem}

\begin{proof}
\begin{enumerate}
\item The learning update $\Delta \theta$ is a function of the observation and
  reward.
\item From \autoref{cond:hallucination}, the observations for $r_{fab}$
  are identical in both worlds.
\item Therefore, the reward distribution and expected parameter updates (the optimization signal that adjusts model
  parameters during training) are
  identical: $\mathbb{E}[\Delta \theta | w_A] = \mathbb{E}[\Delta \theta |
  w_B]$.
\item The optimizer receives no signal to distinguish the cases.
\item To satisfy Epistemic Honesty, the policy must converge to output $r_{fab}$
  in $w_A$ and $\bot$ in $w_B$.
\item However, since the training signal is identical, the policy must converge
  to the same behavior in both cases (either rewarding $r_{fab}$ in both, or
  penalizing it in both).
\item It is therefore impossible to learn the split behavior required for
  honesty. \qedhere
\end{enumerate}
\end{proof}

\begin{remark}[Empirical Bias]
While the theorem proves only that the behavior must be \emph{identical},
empirical observation shows a strong bias. Supervisors typically prefer
``plausible completion'' over ``abstention'' when truth is unknown. This
effectively breaks the tie in favor of fabrication.
\end{remark}

\subsection{Stacking Judges Cannot Escape}

A natural response to \autoref{thm:learnability} is to propose layered
supervision: if one judge is bounded, stack more judges. Let the second judge
supervise the first, the third supervise the second. Surely sufficient layers
will catch fabrications that slip through?

The Observation Monotonicity Lemma closes this escape.

\begin{lemma}[Observation Monotonicity]\label{lem:monotone}
Consider any stack of supervisors $(S_1,\dots,S_k)$ operating under the
text-only observation model. Let $Obs_i(q,r,w)$ denote the observation exported
to $S_i$. If each supervisor can only act on the exported interface emitted by
the previous layer, then there exist deterministic functions $f_i$ such that
$Obs_{i+1}(q,r,w) = f_i(Obs_i(q,r,w))$ for all $i < k$. Consequently, later
judges see no strictly finer epistemic information than earlier ones, regardless
of whether their judgments are deterministic or probabilistic.
\end{lemma}

\begin{proof}
By definition, each supervisor consumes only the serialized response (and any
metadata permitted by the text-only observation model) and produces a finite
judgment that forms the entire input to the next layer. Because no layer can
query the internal world state directly or issue additional interactive probes,
the exported observation available to $S_{i+1}$ must be a deterministic function
of $Obs_i$. Therefore, the informational content across the stack is
monotonically non-increasing, and probabilistic scoring cannot introduce
distinctions absent from $Obs_i$.
\end{proof}

\begin{corollary}[Induction on Layered Judges]\label{cor:layered}
For $k$-layered supervisors $(S_1$ supervising $\pi$, $S_2$ supervising $S_1$,
$\dots$, $S_k$ supervising $S_{k-1})$:
\begin{itemize}
\item \textbf{Base case.} For $k=1$, \autoref{thm:learnability} holds (a
  single bounded supervisor $S$ fails).
\item \textbf{Induction step.} Assume the statement holds for $k=m$. By
  \autoref{lem:monotone}, $S_{m+1}$ receives observations that are
  deterministic functions of $Obs_m$, and thus inherits both the verification
  budget $B$ and the indistinguishability of plausible fabrications. No learning
  or judging signal distinguishes $w_A$ from $w_B$, even if $S_{m+1}$ produces
  probabilistic confidence scores. Therefore the hallucination deadlock persists
  for $k=m+1$.
\end{itemize}
By induction, no finite $k$ resolves the deadlock.
\end{corollary}

\begin{remark}[Scope of Layering]
The inductive result applies only to verification stacks whose layers observe
and transform the same exported response channel. Architectures that surface
additional epistemic state between layers, such as retrieval traces, entity
resolution tables, or cryptographically bound provenance, leave the text-only
observation model and therefore fall outside the impossibility class.
\end{remark}

\subsection{Verification Cost}

The bounded-supervisor assumption deserves scrutiny. Why should verification
cost grow with response complexity? The Composition Graph formalizes this
intuition.

\begin{definition}[Composition Graph]
Given a response $r$, define its composition graph $G(r) = (V, E)$ where each
node $v \in V$ represents an atomic subclaim, named entity, or intermediate
entailment, and each edge $(u, v) \in E$ encodes a dependency that must hold for
$r$ to remain coherent. Edges may represent entailment/support relations,
co-reference constraints, causal ordering, or temporal consistency requirements.
\end{definition}

\begin{lemma}[Superlinear Verification Cost]\label{lem:superlinear}
Under a text-only observation model, reconstructing $G(r)$ from the exported
string requires the supervisor to infer both the nodes and their dependencies
directly from $r$. Even if parsing each node is constant-time, checking
cross-claim consistency requires touching every edge, so the verification cost
is $\Omega(|E|)$. For natural-language responses where each subclaim
participates in multiple constraints (co-reference, causality, temporal order),
$|E|$ grows superlinearly in $|V|$. Consequently, bounded supervisors must
either abandon exhaustive checks or resort to sampling once the number of
subclaims exceeds a modest threshold.
\end{lemma}

\subsection{Responsibility Concentration}

\begin{corollary}[Responsibility Concentration]\label{cor:responsibility}
Under a text-only observation model, epistemic honesty cannot be externally
verified by users or downstream auditors. Responsibility for epistemic honesty
therefore resides solely with the system owner, who retains access to the
internal causal state.
\end{corollary}

External attestation mechanisms, such as cryptographic hash chains, trusted execution
environments, or tamper-evident logs, do not contradict this corollary; rather,
they reinforce it. Such mechanisms only create verifiable guarantees when the
system owner chooses to bind and export epistemic traces, so the authority to
make honesty attestable still rests with the owner.

\subsection{Summary}

The impossibility is now fully characterized. Under text-only observation models
with bounded supervision:
\begin{itemize}
\item The policy cannot satisfy epistemic honesty across ambiguous queries
  (\autoref{thm:representational})
\item The learning algorithm cannot converge to honest behavior
  (\autoref{thm:learnability})
\item Stacking supervisors cannot escape the constraint
  (\autoref{lem:monotone}, \autoref{cor:layered})
\item Verification cost grows superlinearly with response complexity
  (\autoref{lem:superlinear})
\item Responsibility for honesty concentrates with the system owner
  (\autoref{cor:responsibility})
\end{itemize}

These are not limitations to be overcome by better training or more careful
prompting. They are structural properties of the architectural class.

\fakepara{Relation to known results.}
These results formalize what practitioners have observed informally: models
fabricate, and text-only inspection cannot reliably catch it. The formalization
matters because it delineates the boundary precisely. Xu et
al.~\cite{xu2024hallucination} prove that hallucination is computationally
inevitable for LLMs over the class of computable functions, which are a limit on what the
model can \emph{produce}. Our result is orthogonal: a limit on what a supervisor
can \emph{verify}. The RLHF literature documents that text-channel reward is
gameable; we prove that gameability is architectural, not a training failure.
The Fischer-Lynch-Paterson theorem~\cite{flp} showed that consensus is
impossible without timing assumptions, not that distributed systems are hopeless.
Similarly, our result does not argue against verification. It identifies the
precise interface assumption that must change, and the next section constructs
the escape.

We provide two complementary TLA+ specifications that model these constraints
and demonstrate the escape.

The first specification (\texttt{EpistemicImpossibility.tla}) models the
text-only regime. The system maintains a \texttt{vector\_clock} that perfectly
tracks whether each response is grounded in training data or noise. But the
\texttt{Linearize} action exports only text, so the clock is discarded. The
bounded judge (\texttt{JudgeVerify}) can detect obvious lies but not plausible
ones. TLC model-checking finds a counterexample: a plausible lie passes
\texttt{JudgeVerify} while \texttt{IsHonest} returns false. The
\texttt{Verifiability} invariant fails.

The second specification (\texttt{epistemic\_tensor.tla}) models the tensor
interface escape. Now \texttt{ExportTensor} emits a triple: text, provenance,
and topology. The tensor-aware judge (\texttt{TensorVerify}) checks all three
channels. TLC confirms that the \texttt{Verifiability} invariant holds: no lie
escapes the Level~2 judge. The escape is constructive.

Together, the specifications demonstrate both the impossibility and its
resolution. The first shows why text-only observation fails; the second shows
that exporting epistemic metadata restores verifiability. Both specifications
and their TLC traces are available in supplementary materials.

The next section characterizes what it takes to escape this class.
